%LTeX: language=fr-FR
\documentclass[12pt,a4paper]{article}

% ---------- Langue et encodage ----------
\usepackage[french]{babel}
%\usepackage[T1]{fontenc}
%\usepackage[utf8]{inputenc}

% ---------- Mise en page ----------
\usepackage[margin=2.2cm]{geometry}
\usepackage{array,booktabs,tabularx}
\usepackage{amsmath,amssymb}
\usepackage{xcolor}
\usepackage{tcolorbox}
\tcbuselibrary{skins,breakable}
\usepackage{enumitem}
\usepackage{graphicx}
\usepackage{tikz}
\usepackage{pgfplots}
\pgfplotsset{compat=1.18}
\usepackage{fancyhdr}
\usepackage{makecell}
\usepackage{pgffor}
\usepackage{currfile}
\usepackage{lastpage}

\graphicspath{ {../../Images/} }
\input{\currfiledir ../def.tex}
\def \Document {Activité 1}


% ---------- Couleurs ----------
\definecolor{bleuF}{RGB}{47,87,163}   
\definecolor{orH}{RGB}{242,170,60}    
\definecolor{cadre}{RGB}{47,87,163}

% ---------- Boîtes ----------
\newtcolorbox{objectif}{
  colback=bleuF!6, colframe=cadre, boxrule=0.6pt, arc=2pt,
  left=6pt, right=6pt, top=4pt, bottom=4pt,
  title={\textbf{Objectif}}, fonttitle=\bfseries,
  coltitle=cadre, colbacktitle=bleuF!12
}
\newtcolorbox{bilan}{
  colback=orH!8, colframe=cadre, boxrule=0.6pt, arc=2pt,
  left=6pt, right=6pt, top=4pt, bottom=4pt,
  title={\textbf{Bilan}}, fonttitle=\bfseries,
  coltitle=cadre, colbacktitle=orH!20,
  breakable
}
\newtcolorbox{correctionbox}{
  colback=green!5, colframe=green!40!black, boxrule=0.6pt, arc=2pt,
  left=6pt, right=6pt, top=4pt, bottom=4pt,
  breakable
}

% \checkmark symbol fallback (amssymb provides \checkmark)

% ---------- Ligne de réponse pointillée ----------
\newcommand{\reponse}[1]{%
  \par\vspace{2mm}
  \foreach \i in {1,...,#1}{\rule{0pt}{0pt}\dotfill\\[3mm]}%
}

\begin{document}


\pagestyle{fancy}

\fancyhf{} % clear existing header/footer entries
\fancyhead[R]{ {\small \Niveau}}
\fancyhead[L]{ {\small \ChapitreCourt}}
\fancyhead[C]{ {\small \Document}}
%\fancyhead[R]{\large \textbf{\textbf{\Document\ (B) -- 1h}}}
\fancyfoot[L]{ A. Fernandez }
\fancyfoot[R]{{\small \thepage\  / \pageref{LastPage}}}


% ---------- Dessins : dauphin et poisson ----------
% \dauphin{x}{y}{angle}{échelle} : dauphin stylisé, le point (x ; y) est à la base de la queue
\newcommand{\dauphin}[4]{%
  \begin{scope}[shift={(#1,#2)}, rotate=#3, scale=#4, transform shape]
    \filldraw[fill=gray!55!blue!45, draw=gray!40!blue!70!black, line width=0.4pt, line join=round]
      (0,0) .. controls (0.6,0.55) and (1.7,0.95) .. (2.5,0.5)
      .. controls (2.7,0.38) .. (3.0,0.18) -- (3.05,0.1)
      .. controls (2.85,0.05) .. (2.6,0.02)
      .. controls (1.9,-0.45) and (0.7,-0.35) .. (0,0) -- cycle;
    \filldraw[fill=gray!55!blue!45, draw=gray!40!blue!70!black, line width=0.4pt, line join=round]
      (0.12,0.04) .. controls (-0.25,0.15) .. (-0.55,0.45) .. controls (-0.45,0.12) .. (-0.32,0)
      .. controls (-0.45,-0.12) .. (-0.55,-0.45) .. controls (-0.25,-0.15) .. (0.12,-0.04) -- cycle;
    \filldraw[fill=gray!55!blue!45, draw=gray!40!blue!70!black, line width=0.4pt, line join=round]
      (1.15,0.6) .. controls (1.3,1.0) .. (1.65,1.12) .. controls (1.5,0.85) .. (1.65,0.7) -- cycle;
    \filldraw[fill=gray!55!blue!45, draw=gray!40!blue!70!black, line width=0.4pt, line join=round]
      (1.95,-0.22) .. controls (1.85,-0.5) .. (1.55,-0.65) .. controls (1.62,-0.38) .. (1.6,-0.24) -- cycle;
    \fill[white, opacity=0.7] (0.45,-0.12) .. controls (1.3,-0.33) and (2.1,-0.25) .. (2.6,0.0)
      .. controls (2.1,-0.12) and (1.3,-0.15) .. (0.45,-0.12);
    \fill[black] (2.4,0.3) circle (0.05);
  \end{scope}}
% \poisson{x}{y}
\newcommand{\poisson}[2]{%
  \begin{scope}[shift={(#1,#2)}, scale=0.35]
    \fill[orange!80] (0,0) ellipse (0.45 and 0.22);
    \fill[orange!80] (-0.4,0) -- (-0.75,0.22) -- (-0.75,-0.22) -- cycle;
    \fill[black] (0.25,0.05) circle (0.04);
  \end{scope}}

\setlength{\parindent}{0pt}
\newcommand{\trou}[1]{\underline{\hspace{#1}}}

\begin{objectif}
Découvrir une fonction du second degré donnée sous forme factorisée : ses racines, sa courbe, son axe de symétrie et son sommet.
\end{objectif}

\section*{Le saut du dauphin}

Dans un bassin, un dauphin aperçoit un poisson lancé par son soigneur. À l'instant $t = 0$, il sort de l'eau. Il attrape le poisson $2$ secondes plus tard, à une hauteur de $2{,}4$ mètres, puis replonge dans l'eau au bout de $6$ secondes.

\begin{center}
\begin{tikzpicture}[x=1.45cm, y=0.9cm]
  % bassin
  \fill[blue!12] (-0.6,0) rectangle (6.8,-0.9);
  \draw[blue!60, thick] (-0.6,0) -- (6.8,0);
  \node[blue!70!black, font=\small] at (4.2,-0.5) {surface de l'eau : hauteur $0$};
  % trajectoire
  \draw[thick, dashed, blue!70!black, domain=0:6, samples=80] plot (\x,{-0.3*\x*(\x-6)});
  % poisson attrapé à t = 2
  \draw[dotted, thick] (2,0) -- (2,2.4);
  \node[font=\small, right] at (2,1.2) {$2{,}4$ m};
  \poisson{2.05}{2.52}
  \draw[->, gray] (0.9,2.75) node[left, black, font=\small] {poisson lancé par le soigneur} .. controls (1.4,2.85) .. (1.82,2.6);
  % dauphins
  \dauphin{-0.5}{-0.6}{48}{0.62}
  % instants
  \foreach \t/\lab in {0/{$t = 0$}, 2/{$t = 2$ s}, 6/{$t = 6$ s}} {
    \fill (\t,0) circle (1.6pt);
    \node[font=\small, below] at (\t,-0.9) {\lab};
    \draw[gray, very thin] (\t,0) -- (\t,-0.9);
  }
\end{tikzpicture}
\end{center}

On note $f(t)$ la hauteur du dauphin, en mètres, au-dessus de l'eau à l'instant $t$, en secondes, pour $t$ entre $0$ et $6$. La surface de l'eau correspond à la hauteur $0$.

On modélise la trajectoire par une fonction de la forme
\[ f(t) = a\,(t - t_1)(t - t_2), \]
où $a$, $t_1$ et $t_2$ sont des nombres à déterminer, avec $a \neq 0$.

\begin{enumerate}[label=\textbf{\arabic*.}, leftmargin=1.3em]
\item À quels instants le dauphin est-il au niveau de l'eau~? Que vaut alors $f(t)$~?
\reponse{1}

\item Calculer $f(t_1)$ et $f(t_2)$. En déduire les valeurs de $t_1$ et $t_2$, puis compléter :
\[ f(t) = a\,(t - \trou{1cm})(t - \trou{1cm}) = a\,t\,(t - \trou{1cm}). \]
\reponse{2}

\item Le dauphin est à $2{,}4$ m de hauteur à l'instant $t = 2$ : on a donc $f(2) = 2{,}4$. En déduire la valeur de $a$.
\reponse{2}

\item On admet que $f(t) = -0{,}3\,t\,(t - 6)$. Compléter le tableau de valeurs.
\begin{center}
\renewcommand{\arraystretch}{1.8}
\begin{tabular}{|c|*{7}{>{\centering\arraybackslash}p{1.2cm}|}}
\hline
$t$ & $0$ & $1$ & $2$ & $3$ & $4$ & $5$ & $6$ \\
\hline
$f(t)$ & & & & & & & \\
\hline
\end{tabular}
\end{center}

\item Placer les points du tableau dans le repère, puis tracer la courbe de $f$. Cette courbe s'appelle une \textbf{parabole}.
\begin{center}
\begin{tikzpicture}
\begin{axis}[
  width=0.75\textwidth, height=6.5cm,
  xmin=0, xmax=6.5, ymin=0, ymax=3.2,
  xtick={0,...,6}, ytick={0,0.5,...,3}, minor x tick num=1, minor y tick num=1,
  grid=both, major grid style={gray!45}, minor grid style={gray!20},
  axis lines=left, xlabel={$t$ (s)}, ylabel={hauteur (m)},
  ylabel style={font=\small}, xlabel style={at={(1,0)}, anchor=west},
  yticklabel style={/pgf/number format/use comma},
  ]
\end{axis}
\end{tikzpicture}
\end{center}

\item Comparer $f(1)$ et $f(5)$, puis $f(2)$ et $f(4)$. Que remarque-t-on~? Quelle droite est un axe de symétrie de la courbe~? Comment l'obtenir à partir de $t_1$ et $t_2$~?
\reponse{2}

\item Quelle est la hauteur maximale atteinte par le dauphin~? À quel instant~?
\reponse{2}

\item Développer l'expression $f(t) = -0{,}3\,t\,(t - 6)$.
\reponse{2}
\end{enumerate}

\begin{bilan}
\begin{itemize}[leftmargin=1.2em]
\item Une fonction de la forme $f(x) = a\,(x - x_1)(x - x_2)$, avec $a \neq 0$, est une fonction du second degré écrite sous \textbf{forme factorisée}. Sa courbe est une \trou{2.5cm}.
\item Les nombres $x_1$ et $x_2$ sont les \textbf{racines} de $f$ : ce sont les solutions de l'équation $f(x) = \trou{1cm}$.
\item La droite d'équation $x = \trou{2cm}$, à mi-chemin entre les deux racines, est un \textbf{axe de symétrie} de la courbe.
\item Le point de la courbe situé sur cet axe est le \textbf{sommet}. Ici $a < 0$ : la parabole est tournée vers le \trou{1.5cm} et le sommet est le point le plus \trou{1.5cm}.
\end{itemize}
\end{bilan}

\end{document}
